\chapter{Patterns: command queues}
\label{ch:commands}

\begin{aside}
Some actions are commands: user-invoked operations with a
defined name and effect. Storing them in a queue enables
undo, replay, and macros.
\end{aside}

\section{The command pattern}

\begin{lstlisting}
struct Command {
    std::string name;
    std::function<void()> execute;
    std::function<void()> undo;
};
\end{lstlisting}

Each user action constructs a Command and submits.

\section{Dispatching}

A command dispatcher executes and records:

\begin{lstlisting}
class Dispatcher {
    std::vector<Command> history;
public:
    void dispatch(Command c) {
        c.execute();
        history.push_back(c);
    }
    void undo() {
        if (history.empty()) return;
        auto c = history.back();
        history.pop_back();
        c.undo();
    }
};
\end{lstlisting}

\section{Naming}

The name field enables:

\begin{itemize}[leftmargin=*]
  \item Command palette (search by name).
  \item Macro recording (save a sequence).
  \item Telemetry (count which commands are popular).
\end{itemize}

\section{Macros}

A macro is a saved command sequence. Replay by dispatching
each in order.

\section{Recording}

A toggle records every dispatched command. On stop, save
to a macro.

\section{Recap}

\begin{recap}
\begin{itemize}[leftmargin=*]
  \item Commands have name, execute, undo.
  \item Dispatcher records history.
  \item Names enable palette, macros, telemetry.
\end{itemize}
\end{recap}

\section*{Workshop}

\begin{enumerate}[leftmargin=*]
  \item Convert one button to a command.
  \item Wire undo via command history.
  \item Implement macro record and replay.
\end{enumerate}
